\section{Consequences and remaining questions}
The reduction makes a linear-time dyadic bit transposer sufficient
for canonical word transposition. Its contrapositive says that failure of linear-time canonical word
transposition implies failure of linear-time dyadic bit transposition.
The reverse transfer would require $\WORD\Rightarrow\BIT$, which
remains open. No general quantitative lower-bound transfer is asserted.

The multiscale argument isolates a common obstruction: a short composition of
whole-entry operations and simple rearrangements has limited fine-scale
spread unless it traverses a block whose scale is exceptional. Disjointness
bounds the total volume of these exceptional blocks across sufficiently
separated scales. To cover synchronous scans, path counting must be supplemented by a bound
on information carried between addresses. For the synchronous scan grammar, one shared record supplies the
exceptional values and scan boundary states, after which acyclic decoding
charges the missing named input bits. The finite bound quantifies this cost:
a helper of $h$ bits can support at most $(D+1)h$ demands whose own bins
are absent from ancestor sets of size at most $D$. The count is derived even when the intermediate bank encodings are nonlinear.

The state restriction matters. The theorem permits arbitrary nonlinear
transition tables, but fixes the number of banks, layers, and state bits while
the dimension grows. It does not permit an unbounded workspace to travel
through a scan without being charged. Nor does it cover arbitrary asynchronous
head schedules, input-dependent block layouts, or expanding banks. Extending
the barrier to any of these features requires an additional proof.

The construction program now has a proved sufficient target: a uniform
linear-time dyadic bit transposer supplies the canonical word transposer
through \cref{thm:bit-word}. Alternatively, a reverse compiler must use
operations beyond the excluded grammars; allowing such operations is only
a candidate route, not evidence of a construction. In either case, every
movement, state update, temporary copy and reset must be charged in the
native model. For the impossibility program, the next target is an accounting
principle that continues to bound the shared information when local state
or access schedules grow with the input. Neither a
restricted negative result nor a short algebraic factorization alone settles
the original transposition question.
